\documentclass[11pt,a4paper]{article}
\usepackage{DDTA_DERMATRIAGE_GUIDED_REVIEW_STYLE_R1}

\hypersetup{
  pdftitle={DDTA R25 DermaTriage - Guided Documentation Review R3},
  pdfauthor={Documentation-Driven Threat Analysis research project}
}

\pagestyle{fancy}
\fancyhf{}
\fancyhead[L]{\sffamily\small\color{DDTAGray}DDTA DermaTriage Guided Review}
\fancyhead[R]{\sffamily\small\color{DDTAGray}R25 - WORKING RECORD R3}
\fancyfoot[L]{\sffamily\scriptsize\color{DDTAGray}Documentation reconstruction before Base Analysis}
\fancyfoot[R]{\sffamily\small\color{DDTAGray}\thepage}
\renewcommand{\headrulewidth}{0.4pt}
\renewcommand{\footrulewidth}{0pt}

\begin{document}
\selectlanguage{italian}

% -----------------------------------------------------------------------------
% TITLE
% -----------------------------------------------------------------------------
\begin{titlepage}
\thispagestyle{empty}
\vspace*{8mm}

{\sffamily\bfseries\color{DDTANavy}\fontsize{15}{18}\selectfont Documentation-Driven Threat Analysis}\par
\vspace{5mm}
{\sffamily\bfseries\color{DDTANavy}\fontsize{29}{33}\selectfont DermaTriage Guided Documentation Review}\par
\vspace{2mm}
{\sffamily\color{DDTABlue}\fontsize{18}{22}\selectfont Working Record R3 - chiusura cumulativa della Macro Project Map}\par

\vspace{10mm}
\begin{tcolorbox}[enhanced,arc=3pt,boxrule=0pt,colback=DDTANavy,left=12pt,right=12pt,top=11pt,bottom=11pt]
{\color{white}\sffamily\bfseries\large Registro cumulativo della validazione holdout DermaTriage}\par
\vspace{2mm}
{\color{white!88}\small Conserva lo stato cumulativo R2 e aggiunge il reopen controllato di A2, la stabilizzazione della Macro Project Map con MR-02, MR-03 e MR-04, il rifiuto controllato del candidate diagnostico e il gate cumulativo D1. Solo dopo il documentation gate potra' essere derivata la Base Analysis.}
\end{tcolorbox}

\vspace{8mm}
\begin{tabularx}{\textwidth}{L{42mm}Y}
\toprule
\textbf{Stato} & R25 WORKING REVIEW RECORD - DOCUMENTATION PHASE \\
\textbf{Repository} & \texttt{nballestriero/documentation-driven-threat-analysis} \\
\textbf{Baseline R1 originario} & \texttt{21575a18ccf2feab62b27e3b046fadcd01c90135} \\
\textbf{Baseline repository corrente} & \texttt{4dd52e02cc962aee49d2c4593ff851802338698d} \\
\textbf{Metodo} & \texttt{DDTA\_DOCUMENTATION\_BA\_AUTHORING\_GUIDE\_R4} \\
\textbf{Caso holdout} & DermaTriage - ricostruzione da documentazione esistente \\
\textbf{Base Analysis} & BLOCCATA fino al superamento del documentation gate \\
\textbf{Data} & 4 settembre 2026 \\
\bottomrule
\end{tabularx}

\vfill
\begin{ddtaprinciple}[title={Ruolo di questo artifact}]
Questo documento non e' la documentazione finale DermaTriage e non e' una Base Analysis. E' il registro guidato della ricostruzione: conserva source evidence, candidate meaning, gate applicati, esiti, gap e osservazioni metodologiche. Al termine della review dovra' rendere trasparente non soltanto \emph{cosa} e' stato scritto, ma \emph{perche'} ogni elemento DDTA e' stato accettato, modificato, fermato o lasciato non specificato.
\end{ddtaprinciple}
\end{titlepage}

\tableofcontents
\clearpage

% -----------------------------------------------------------------------------
\section{Obiettivo della sessione e disciplina di lavoro}

L'obiettivo concordato e' arrivare alla scrittura della documentazione DermaTriage in modo guidato, usando un progetto reale per rendere osservabile la metodologia DDTA. La review deve inoltre produrre evidence sufficiente per concludere, a valle del lavoro, quali siano i pregi e i difetti di DDTA e in che modo il metodo contribuisca a migliorare la documentazione di progetto.

\begin{ddtacore}[title={Sequenza autorizzata}]
\begin{lstlisting}
repository baseline
    -> methodology reading
    -> original DermaTriage package + SHA-256
    -> authority gate
    -> project problem framing
    -> MacroRequirement
    -> semantic sufficiency
    -> Decision
    -> FunctionalRequirement
    -> Requirement split
    -> Specialized / Security requirements
    -> cross-MR and terminology review
    -> documentation completeness / promotion gate
    -> ONLY THEN Base Analysis
\end{lstlisting}
\end{ddtacore}

La regola collaborativa della sessione e' una sola: nessuna classificazione dell'intero progetto viene eseguita silenziosamente. Ogni elemento viene discusso e chiuso prima di avanzare.

\begin{lstlisting}
source evidence
    -> candidate project meaning
    -> relevant R4 gates
    -> PASS / REWORK / NOT SPECIFIED / CONFLICTING / METHOD PRESSURE
    -> next element only after disposition
\end{lstlisting}

\begin{ddtaanti}[title={Anti-contamination rule}]
Durante questa fase non si usa Base Analysis per decidere project meaning, struttura dei documenti o livello di decomposizione. I vecchi candidate DDTA di DermaTriage possono essere usati in futuro come evidence comparativa, ma non come autorita' sul progetto.
\end{ddtaanti}

% -----------------------------------------------------------------------------
\section{Authority metodologica e repository baseline}

Prima della source review sono stati letti, nell'ordine concordato:
\begin{enumerate}
  \item \texttt{DDTA\_R25\_DOCUMENTATION\_METHOD\_STABILIZATION\_CHECKPOINT\_R1.md};
  \item \texttt{DDTA\_R25\_HOLDOUT\_VALIDATION\_WORK\_PLAN\_R1.md};
  \item \texttt{DDTA\_DOCUMENTATION\_BA\_AUTHORING\_GUIDE\_R4.tex};
  \item \texttt{DDTA\_R25\_DERMATRIAGE\_DOCUMENTATION\_REVIEW\_CONTINUATION\_HANDOFF\_R1.md}.
\end{enumerate}

Il branch \texttt{master} e' stato verificato sul commit:
\begin{lstlisting}
21575a18ccf2feab62b27e3b046fadcd01c90135
\end{lstlisting}

Questo commit e' il continuation baseline originario della sessione R1. Dopo i checkpoint di preservazione degli artifact di review e la chiusura della Decision family di MR-01, il baseline repository corrente per R3 e' \texttt{4dd52e02cc962aee49d2c4593ff851802338698d}. Il metodo R4 stabilizzato e' il forward authoring protocol; gli esperimenti R5 storici non sono forward authority.

\begin{ddtacheck}[title={Repository authority gate}]
\textbf{PASS.} La sessione parte dal baseline esplicitamente richiesto e non da una ricostruzione cronologica o da un artifact piu' recente scelto per recency.
\end{ddtacheck}

% -----------------------------------------------------------------------------
\section{A0 - Original source package e authority gate}

\subsection{Package pinned}

Il work plan e l'handoff identificano come unica source authority primaria per il significato DermaTriage il package:

\begin{lstlisting}
DermaTriage-Docs-20260830T152637Z-1-001.zip
Expected SHA-256:
E9ED2C507BEFB95F54A52084687CD1E8798863AE81CF69D09568864D8CBF280E
\end{lstlisting}

Il package originale e' stato riaperto e l'hash e' stato ricalcolato sui byte materializzati:

\begin{lstlisting}
Computed SHA-256:
E9ED2C507BEFB95F54A52084687CD1E8798863AE81CF69D09568864D8CBF280E
\end{lstlisting}

\begin{ddtacore}[title={A0 disposition}]
\textbf{PASS.} Nome e SHA-256 coincidono esattamente con il package pinned. Non viene sostituito con archivi simili o con working package prodotti durante precedenti passaggi DDTA.
\end{ddtacore}

\subsection{Source inventory}

L'archivio contiene sette file:

\begin{tabularx}{\textwidth}{L{59mm}Y}
\toprule
\textbf{Source artifact} & \textbf{Ruolo iniziale nella review} \\
\midrule
\nolinkurl{OR2_Architecture_Document.pdf} & Descrizione di purpose, pipeline, integrazioni, API e lifecycle operativo. \\
\nolinkurl{OR2_Model_Test_Report.pdf} & Evidence di modello, metriche, threshold e risultati di test. \\
\nolinkurl{OR3_Dataset_Metadata_Catalog.pdf} & Evidence su dataset, schema e uso dichiarato. \\
\nolinkurl{OR4_Training_Environment_Config.pdf} & Realization/configuration evidence per training e runtime. \\
\nolinkurl{OR4_Training_Cycles_Report.pdf} & Evidence di training e valutazione. \\
\nolinkurl{OR5_Test_Environment_Setup.pdf} & Test/integration evidence, endpoint e pass condition. \\
\nolinkurl{efficientnet_b4.pth} & Binary model evidence; non normative project meaning per default. \\
\bottomrule
\end{tabularx}

\begin{ddtaprinciple}[title={Authority non significa promozione automatica}]
Essere parte del package autoritativo significa che il fatto source-supported deve essere considerato; non significa che ogni fatto tecnico diventi automaticamente \MR, \DEC{} o \REQ. Test, configuration, runtime state e current realization devono essere classificati prima di ricevere status normativo.
\end{ddtaprinciple}

\subsection{A0 gate summary}

\begin{tabularx}{\textwidth}{L{55mm}L{27mm}Y}
\toprule
\textbf{Check} & \textbf{Disposition} & \textbf{Nota} \\
\midrule
Repository baseline & PASS & Commit esatto verificato. \\
Package name & PASS & Nome esatto pinned. \\
Package SHA-256 & PASS & Hash calcolato = hash atteso. \\
Original source as project authority & PASS & Fonte primaria stabilita. \\
Previous DDTA candidates as authority & EXCLUDED & Solo future comparison evidence. \\
Base Analysis as authority & EXCLUDED & Documentation phase ancora aperta. \\
\bottomrule
\end{tabularx}

% -----------------------------------------------------------------------------
\section{A1 - Project problem framing}

\subsection{Source evidence osservata}

La fonte architetturale dichiara esplicitamente:

\begin{ddtaexample}[title={OR2 Architecture Document - purpose}]
\textbf{Purpose:} ``AI-assisted dermatology triage for early skin cancer detection.''
\end{ddtaexample}

La stessa fonte descrive una pipeline che arricchisce il contesto clinico prima della ``final triage decision'', una adaptation layer che assegna una priorita' P1--P4, specialista e SLA, e una successiva doctor validation/correction.

Per il framing, i nomi delle tecnologie e le scelte correnti di soluzione vengono temporaneamente rimossi: EfficientNet, Qwen2-VL, ChromaDB, BioMistral, B4, endpoint, threshold e storage non sono necessari per esprimere il problema generale.

\subsection{Candidate project problem framing}

\begin{ddtacore}[title={Candidate framing accettato nella sessione}]
DermaTriage affronta il problema di supportare il triage precoce di casi dermatologici relativi a lesioni cutanee potenzialmente oncologiche, utilizzando le informazioni disponibili sul caso per determinare una priorita' di presa in carico e supportare l'instradamento verso l'assistenza specialistica appropriata.
\end{ddtacore}

\subsection{Gate D0}

\begin{tabularx}{\textwidth}{L{54mm}L{23mm}Y}
\toprule
\textbf{Domanda} & \textbf{Esito} & \textbf{Motivazione} \\
\midrule
Usa vocabolario di problema/dominio? & PASS & Triage dermatologico, caso, priorita', presa in carico. \\
Resta valido senza la soluzione corrente? & PASS & Nessun modello, provider, API o database e' necessario. \\
Boundary significativo comprensibile? & PASS con gap esplicito & Il triage e' chiaro; l'autorita' diagnostica definitiva non e' stabilita. \\
E' leggibile da un domain reader? & PASS & Non richiede conoscenza dell'architettura corrente. \\
\bottomrule
\end{tabularx}

% -----------------------------------------------------------------------------
\section{Prima semantic-sufficiency question: triage vs diagnosi}

Durante il framing emerge la prima distinzione materiale che il metodo impedisce di nascondere dietro la terminologia della soluzione.

La source documentation usa contemporaneamente:
\begin{itemize}
  \item il purpose di \emph{dermatology triage};
  \item un output Stage 4 chiamato \texttt{pathology};
  \item endpoint e storage denominati \texttt{diagnose}/\texttt{diagnoses};
  \item una successiva doctor validation/correction.
\end{itemize}

\begin{ddtacheck}[title={Smallest critical proposition}]
DermaTriage ha la responsabilita' e l'autorita' di determinare una diagnosi clinica definitiva, oppure la \texttt{predicted\_pathology} e' informazione analitica di supporto al triage e alla successiva review medica?
\end{ddtacheck}

\begin{ddtaopen}[title={Disposition corrente: NOT SPECIFIED}]
La source documentation esaminata non stabilisce in modo sufficiente che DermaTriage abbia autorita' per una diagnosi clinica definitiva. La distinzione non viene completata per inferenza. La \texttt{predicted\_pathology} resta tuttavia un fatto source-supported materiale che dovra' essere conservato e classificato al semantic owner appropriato.
\end{ddtaopen}

\subsection{Diagnosi del problema}

\begin{tabularx}{\textwidth}{L{49mm}L{28mm}Y}
\toprule
\textbf{Candidate diagnosis} & \textbf{Esito} & \textbf{Motivo} \\
\midrule
Source/documentation gap & YES & La fonte usa termini che permettono letture materialmente diverse. \\
Authoring-guide problem & NON DIMOSTRATO & La guida ha fatto emergere la distinzione. \\
Metamodel pressure & NON DIMOSTRATO & Non esiste ancora un significato governato che il metamodel non riesca a rappresentare. \\
BA issue & NON APPLICABILE & BA non e' iniziata. \\
\bottomrule
\end{tabularx}

\begin{ddtaprinciple}[title={Osservazione metodologica 01}]
In questo caso DDTA migliora la documentazione non aggiungendo dettaglio, ma obbligando a separare un termine tecnico presente nell'implementazione da una responsabilita' clinica che la fonte non governa chiaramente. Il valore osservato e' la visibilita' del gap, non la sua chiusura artificiale.
\end{ddtaprinciple}

% -----------------------------------------------------------------------------
\section{A2 - Primo MacroRequirement}

\subsection{Candidate MR-01}

\begin{ddtacore}[title={MR-01 - Valutazione di triage del caso dermatologico}]
\textbf{Intent.} Determinare, dalle informazioni disponibili sul caso dermatologico, una valutazione di urgenza utile a stabilire la priorita' di presa in carico specialistica.

\textbf{Context.} DermaTriage supporta il triage precoce di lesioni cutanee potenzialmente oncologiche.

\textbf{Stakeholders.} Paziente; professionista sanitario coinvolto nella presa in carico e nella review.

\textbf{Scope.} Valutazione del caso fino alla determinazione dell'urgenza/priorita' di triage.

\textbf{Assumptions / Constraints.} Nessuna necessaria al livello MR in questa fase.

\textbf{dependsOn.} Nessuna relazione macro introdotta in questa fase.
\end{ddtacore}

Il candidate MR non incorpora EfficientNet, quattro stage, B4, P1--P4, threshold, pathology prediction o SLA. L'assenza e' intenzionale: questi significati devono essere classificati nei livelli appropriati invece di essere caricati nella macro-responsabilita'.

\subsection{Gate D1 - MR}

\begin{longtable}{L{49mm}L{21mm}L{76mm}}
\toprule
\textbf{Test R4} & \textbf{Esito} & \textbf{Lettura della sessione} \\
\midrule
\endhead
Una sola macro-responsabilita' & PASS & Governa la valutazione di triage del caso. \\
Problem anchoring & PASS & Deriva direttamente dal framing. \\
Architecture / technology resilience & PASS & Sopravvive alla sostituzione completa della soluzione. \\
Temporal stability & PASS & Non dipende dalla configurazione corrente. \\
No comportamento atomico da FR & PASS & Non prescrive stage, chiamate o operazioni. \\
No solution commitment leakage & PASS & Nessun modello, API o provider. \\
Dependency distinta da containment & PASS & Nessuna dependency inventata. \\
Broad domain readability & PASS & Comprensibile senza conoscere DermaTriage internamente. \\
Split image vs symptom-only a livello MR & NO SPLIT & Al momento sono due path della stessa responsabilita' di triage, non due macro-responsabilita' governate. \\
Macro project map globale & PENDING & Potra' essere chiusa solo dopo l'identificazione dell'insieme MR. \\
\bottomrule
\end{longtable}

\subsection{Semantic sufficiency del MR}

La lettura materialmente critica e' ancora la distinzione:

\begin{lstlisting}
triage responsibility
        !=
definitive diagnostic responsibility
\end{lstlisting}

La prima e' \texttt{AFFIRMED}; la seconda resta \NS. Il MR e' scritto in modo da non assorbire la seconda.

Un'ulteriore distinzione source-supported -- urgenza analitica HIGH/MEDIUM/LOW contro priorita' operativa P1--P4 con SLA -- e' materiale, ma non richiede uno split dell'MR. Deve essere conservata nella review downstream.

\begin{ddtacore}[title={MR-01 disposition}]
\textbf{PASS.} Il MR e' semanticamente sufficiente per il proprio livello, con il gap sull'autorita' diagnostica preservato esplicitamente.
\end{ddtacore}

\subsection{STOP test}

\begin{ddtacheck}[title={Serve ulteriore significato governato nello scope corrente?}]
\textbf{YES.} La fonte governa ulteriori distinzioni materiali: path con immagine, fallback symptom-only, pipeline analitica, trasformazione dell'urgenza in priorita' operativa, B4 boundary, clinician review/correction e lifecycle di apprendimento. Questi significati devono essere esaminati uno per volta senza forzare in anticipo il loro owning layer.
\end{ddtacheck}

Pertanto il branch non si ferma a MR-01 e la prossima azione autorizzata e' la review della prima candidate \DEC{} sotto MR-01.

% -----------------------------------------------------------------------------
\section{Checkpoint storico R1 - stato dopo il primo MR}
\begin{ddtaopen}[title={Snapshot storico preservato}]
La figura e la tabella di questa sezione fotografano intenzionalmente lo stato di R1, nel quale la first Decision era ancora NEXT. Sono preservate per mostrare l'evoluzione e sono \textbf{superseded} dallo stato R2 riportato nelle sezioni A3/A4 e nel Session snapshot finale.
\end{ddtaopen}

\begin{center}
\begin{tikzpicture}[
  node distance=7mm and 12mm,
  every node/.style={font=\sffamily\small,align=center},
  done/.style={draw=DDTAGreen,rounded corners=2pt,very thick,fill=DDTALightGreen,minimum width=50mm,minimum height=9mm},
  now/.style={draw=DDTAAmber,rounded corners=2pt,very thick,fill=DDTALightAmber,minimum width=50mm,minimum height=9mm},
  later/.style={draw=DDTAGray,rounded corners=2pt,thick,fill=DDTALightGray,minimum width=50mm,minimum height=9mm},
  blocked/.style={draw=DDTARed,rounded corners=2pt,very thick,fill=DDTALightRed,minimum width=50mm,minimum height=9mm},
  arr/.style={-{Latex[length=2.2mm]},thick,draw=DDTANavy}
]
\node[done] (a0) {A0 Authority / source gate\\PASS};
\node[done,below=of a0] (a1) {A1 Problem framing\\PASS + explicit NOT SPECIFIED};
\node[done,below=of a1] (a2) {A2 MR-01 + semantic sufficiency\\PASS};
\node[now,below=of a2] (a4) {NEXT: first Decision candidate\\gate-by-gate review};
\node[later,below=of a4] (rest) {FR / split / SR / SecR / boundaries\\terminology / propagation / completeness};
\node[blocked,below=of rest] (ba) {Base Analysis\\BLOCKED until documentation gate};
\draw[arr] (a0)--(a1);
\draw[arr] (a1)--(a2);
\draw[arr] (a2)--(a4);
\draw[arr] (a4)--(rest);
\draw[arr] (rest)--(ba);
\end{tikzpicture}
\end{center}

\begin{tabularx}{\textwidth}{L{42mm}L{28mm}Y}
\toprule
\textbf{Review object} & \textbf{Status} & \textbf{Disposition} \\
\midrule
Repository baseline & CLOSED & Exact commit verified. \\
Original package & CLOSED & Exact SHA-256 verified. \\
Authority gate & CLOSED & Original source is project authority. \\
Problem framing & CLOSED for current scope & Triage framing accepted; diagnostic authority remains \NS. \\
MR-01 & CLOSED for current scope & PASS; branch continues. \\
First Decision & NEXT & Not yet authored. \\
Documentation gate & OPEN & Most conditions not yet evaluated. \\
Base Analysis & BLOCKED & Explicitly not started. \\
\bottomrule
\end{tabularx}

% -----------------------------------------------------------------------------

% -----------------------------------------------------------------------------
\section{A3/A4 - Decision family review di MR-01}

Questa sezione costituisce l'avanzamento principale di R2. La validation mode rende espliciti, per ogni candidato, \textbf{domanda -> evidence -> risposta -> esito -> gap -> STOP}. Le risposte sono persistite per studiare la metodologia; questo non implica che un futuro editor operativo debba conservarle tutte nel documento governato.

\subsection{Decision family result in sintesi}

\begin{longtable}{L{48mm}L{31mm}L{78mm}}
\toprule
\textbf{Candidate / element} & \textbf{Disposition} & \textbf{Routing stabilizzato} \\
\midrule
\endhead
Continuita' del triage senza immagine & DEC-01 PASS & Decision SELECTABLE sotto MR-01. \\
P-scale come priorita' operativa & DEC-02 REWORK -> PASS & Decision SELECTABLE sotto MR-01; exact mapping separato come future FR candidate. \\
SLA 24h / 48h / 72h / 7d & NOT SPECIFIED & Non promosso a Decision: manca semantic binding temporale e responsibility owner. \\
Specialist output / routing & OUTSIDE MR-01 & Non e' una Decision di MR-01; riapre la macro project map come concern candidato separato. \\
HIGH/MEDIUM/LOW + confidence & EVIDENCE / future FR semantics & Non promosso a Decision indipendente: rappresentazione analitica corrente da riesaminare con il mapping FR. \\
Four-stage AI pipeline e named models & REALIZATION EVIDENCE & Nessuna promozione automatica a Decision; la fonte descrive architettura corrente ma non dimostra un commitment normativo separato necessario a MR-01. \\
B4 persistence / external APIs & CONSUMED SERVICE EVIDENCE & Boundary evidence; non ownership della responsabilita' di triage. \\
Pathology prediction & OUTSIDE / GAP & Material evidence, ma autorita' diagnostica definitiva resta NOT SPECIFIED. \\
Doctor correction / prompt evolution / retraining & OUTSIDE MR-01 & Candidati per macro-responsibility/lifecycle review successiva. \\
\bottomrule
\end{longtable}

% -----------------------------------------------------------------------------
\section{DEC-01 - Continuita' del triage in assenza di immagine}
\closedtag

\subsection{Source evidence}
\textbf{Source:} \texttt{OR2\_Architecture\_Document.pdf}, pipeline/fallback. La fonte descrive un percorso image-based e afferma che, quando l'immagine non e' disponibile, uno scoring symptom-only deriva l'urgenza dai dati sintomatologici disponibili nell'interazione B4.

\subsection{Candidate Decision stabilizzata}
\textbf{Context.} Il caso dermatologico puo' essere disponibile con un'immagine della lesione oppure senza immagine. La documentazione descrive un percorso principale image-based e un fallback symptom-only.

\textbf{Decision.} DermaTriage deve consentire la valutazione di triage anche quando non e' disponibile un'immagine della lesione. In tale condizione, la determinazione dell'urgenza deve basarsi sulle informazioni sintomatologiche disponibili sul caso.

\textbf{Consequences.} Il triage deve supportare almeno le condizioni ``con immagine'' e ``senza immagine''. Non e' autorizzato assumere che il percorso symptom-only produca l'intero insieme di risultati del percorso image-based.

\subsection{Gate review esplicita}
\begin{longtable}{L{50mm}L{24mm}L{83mm}}
\toprule
\textbf{Domanda} & \textbf{Esito} & \textbf{Risposta / evidence} \\
\midrule
\endhead
Esiste un material project commitment? & PASS & L'immagine non e' una precondizione assoluta per poter effettuare il triage. \\
Restringe esattamente MR-01? & PASS & Governa le condizioni di applicabilita' della medesima responsabilita' di triage. \\
E' tautologica rispetto all'MR? & PASS & MR-01 non stabilisce da solo cosa accade in assenza di immagine. \\
Split necessario? & PASS / NO SPLIT & Con immagine e senza immagine sono condizioni della stessa scelta di continuita', non responsabilita' indipendenti. \\
Responsibility boundary? & PASS & DermaTriage possiede il comportamento di triage; B4 fornisce/intermedia dati ed e' capability consumata. \\
Decision vs realization? & PASS & Eliminando chatbot, B4 e scoring corrente rimane la scelta di consentire il triage senza immagine. \\
Decision kind? & SELECTABLE & Una alternativa materialmente distinta sarebbe richiedere l'immagine come precondizione. \\
Semantic sufficiency? & PASS & Il commitment e' stabile; l'equivalenza completa degli output tra i due percorsi non e' specificata. \\
D2 / STOP? & PASS / NO STOP & Restavano altri commitment materiali da classificare sotto MR-01. \\
\bottomrule
\end{longtable}

\begin{ddtaopen}[title={DEC-01 - gap preservati}]
\textbf{NS-01.} Alternative effettivamente considerate e rationale storico non documentati.\par
\textbf{NS-02.} Non e' specificato se il fallback symptom-only produca, oltre all'urgenza, gli stessi outcome del percorso image-based.
\end{ddtaopen}

% -----------------------------------------------------------------------------
\section{DEC-02 - Adozione della P-scale per la priorita' operativa}
\closedtag

\subsection{Source evidence}
\textbf{Source principale:} \texttt{OR2\_Architecture\_Document.pdf}. La fonte distingue una valutazione analitica \texttt{HIGH/MEDIUM/LOW + confidence} da una priorita' operativa \texttt{P1/P2/P3/P4}; la adaptation layer documentata associa inoltre SLA e specialist.

\begin{lstlisting}
HIGH + confidence > 0.85 -> P1 -> 24 hours
HIGH                     -> P2 -> 48 hours
MEDIUM                   -> P3 -> 72 hours
LOW                      -> P4 -> 7 days
\end{lstlisting}

\subsection{Candidate meaning iniziale}
Il candidato spontaneo ``P-scale + exact mapping + SLA + specialist'' viene sottoposto a split gate senza essere assunto come valido.

\subsection{Decision gate}
\begin{ddtacheck}
\textbf{Domanda.} Eliminando i dettagli esecutivi, resta una posizione progettuale che restringe materialmente MR-01?
\end{ddtacheck}
\textbf{Risposta.} SI. Adottare P1-P4 come rappresentazione della priorita' operativa e' una convenzione/strategia sostituibile da convenzioni materialmente diverse.\par
\textbf{Esito: PASS.}

\begin{ddtacheck}
\textbf{Domanda.} MR-01 implica gia' necessariamente P1-P4?
\end{ddtacheck}
\textbf{Risposta.} NO. MR-01 richiede una priorita' di triage ma non prescrive una scala specifica.\par
\textbf{Esito: PASS - non tautologica.}

\subsection{Split gate}
\begin{ddtacheck}
\textbf{Domanda.} P-scale, mapping condizionale, SLA e specialist routing possono cambiare indipendentemente con conseguenze materiali autonome?
\end{ddtacheck}
\textbf{Risposta.} SI. Si puo' mantenere P1-P4 cambiando threshold, SLA o routing specialistico.\par
\textbf{Esito sul candidato iniziale: REWORK.}

\begin{ddtacore}[title={DEC-02 dopo REWORK}]
\textbf{Title:} Adozione della P-scale per la priorita' operativa.\par
\textbf{Context:} la valutazione analitica dell'urgenza deve essere trasformata in una rappresentazione utilizzabile per esprimere la priorita' operativa del triage.\par
\textbf{Decision:} DermaTriage rappresenta la priorita' operativa del triage mediante la scala P1-P4, derivandola dalla valutazione di urgenza del caso.\par
\textbf{Consequences:} il comportamento downstream deve produrre una priorita' appartenente al dominio P1-P4; la regola che seleziona il livello specifico viene governata separatamente.
\end{ddtacore}

\subsection{Responsibility-boundary gate}
\begin{ddtacheck}
\textbf{Domanda.} La determinazione della P-scale appartiene a DermaTriage o a B4/servizio esterno?
\end{ddtacheck}
\textbf{Risposta.} La trasformazione e' descritta nel flusso DermaTriage; B4 riceve/persistisce dati e risultati. DermaTriage possiede la determinazione della priorita' P.\par
\textbf{Esito: PASS.}

\subsection{Decision-vs-realization gate}
\begin{ddtacheck}
\textbf{Domanda.} Rimuovendo \texttt{Adaptation Layer} e il nome della funzione, resta una scelta governata?
\end{ddtacheck}
\textbf{Risposta.} SI. Sostituire il componente non cambia la P-scale; sostituire P1-P4 con una convenzione diversa cambia il significato governato.\par
\textbf{Esito: PASS.}

\subsection{Decision kind}
\textbf{Kind: SELECTABLE.} Alternative materialmente differenti sono possibili, ma la fonte non documenta quali siano state considerate ne' il rationale storico.

\begin{lstlisting}
commitment              = AFFIRMED
alternatives considered = NOT SPECIFIED
historical rationale    = NOT SPECIFIED
\end{lstlisting}

\subsection{Cross-layer gate: Decision o FR?}
\begin{ddtacheck}
\textbf{Domanda.} La truth table \texttt{urgency/confidence -> specific P} appartiene alla Decision?
\end{ddtacheck}
\textbf{Risposta.} NO. La P-scale e' il dominio/convenzione; il mapping e' un obbligo operativo separabile e resta future FR candidate.\par
\textbf{Esito: PASS dopo REWORK.}

\subsection{Semantic sufficiency}
\begin{itemize}
  \item la P-scale derivata dall'urgenza e' \texttt{AFFIRMED};
  \item la confidence e' materialmente rilevante per la regola documentata: \texttt{AFFIRMED};
  \item la semantica temporale dello SLA e' insufficiente: \texttt{NOT SPECIFIED};
  \item l'applicazione della stessa trasformazione P-scale al fallback symptom-only non e' esplicitamente governata: \texttt{NOT SPECIFIED};
  \item la regola che determina lo specialist non e' documentata: \texttt{NOT SPECIFIED}.
\end{itemize}

\textbf{D2 individuale: REWORK -> PASS. STOP individuale al momento della review: NO.}

% -----------------------------------------------------------------------------
\section{Candidate SLA - rifiuto controllato della promozione a Decision}
\holdouttag

\subsection{Source proposition minima}
La fonte associa a \texttt{P1/P2/P3/P4} valori rispettivamente pari a 24h, 48h, 72h e 7 giorni sotto la label \texttt{SLA}. Questa proposizione e' \texttt{AFFIRMED}.

\subsection{Decision gate e split}
Un candidato del tipo ``ogni P-scale possiede un tempo massimo di presa in carico'' sarebbe materialmente significativo, ma introduce gia' una semantica non affermata dalla fonte. Inoltre il principio ``esiste una dimensione temporale'' e i valori numerici specifici possono cambiare indipendentemente.\par
\textbf{Esito: PENDING -> REWORK.}

\subsection{Responsibility-boundary e semantic-sufficiency gate}
\begin{ddtacheck}
\textbf{Domanda.} Che cosa deve accadere entro 24/48/72 ore o 7 giorni, da quale evento decorre il tempo e chi e' responsabile del rispetto?
\end{ddtacheck}

Sono materialmente possibili almeno letture diverse: contatto, appuntamento, presa in carico, handoff, semplice target informativo o altro evento. La fonte non stabilisce trigger temporale, outcome finale, responsibility owner, failure behavior o se il valore sia limite, target o raccomandazione.

\begin{ddtaopen}[title={GAP-DERMA-SLA-01}]
\textbf{Disposition: NOT SPECIFIED - non promuovere a Decision.}\par
La precisione numerica non implica semantic sufficiency. I valori restano source evidence da preservare finche' non sia definito il semantic binding temporale.
\end{ddtaopen}

\textbf{Decision kind: NOT ASSIGNED. STOP candidate: YES.}

% -----------------------------------------------------------------------------
\section{Candidate specialist routing - semantic-owner gate}
\holdouttag

\subsection{Source evidence}
\textbf{OR2 Architecture Document:} la adaptation layer ``assigns P1-P4, specialist and SLA''.\par
\textbf{OR5 Test Environment Setup:} la risposta attesa di \texttt{/analyze} contiene \texttt{P1-P4, specialist, SLA, confidence}; il test full-pipeline attende ``P1-P4 + Italian specialist + SLA time''.

La proposizione minima supportata e': \textbf{la risposta DermaTriage contiene una indicazione specialistica}.\par
\textbf{Evidence classification: AFFIRMED.}

\subsection{Candidate meaning iniziale}
\begin{quote}
DermaTriage determina lo specialista appropriato come parte del triage del caso.
\end{quote}

\subsection{Material commitment gate}
\begin{ddtacheck}
\textbf{Domanda.} Eliminando il nome dell'Adaptation Layer, resta un significato progettuale materiale?
\end{ddtacheck}
\textbf{Risposta.} SI. La presenza di una indicazione specialistica e' osservabile anche nei test e sopravvive ai nomi dei componenti.\par
\textbf{Esito: PASS come candidate project meaning.}

\subsection{Exact-parent / semantic-owner gate}
\begin{ddtacheck}
\textbf{Domanda.} Questo commitment restringe esattamente MR-01 oppure introduce una responsabilita' diversa?
\end{ddtacheck}

\textbf{Evidence.} MR-01 ha scope ``valutazione del caso fino alla determinazione dell'urgenza/priorita' di triage''. Il project framing, separatamente, include il supporto all'instradamento verso assistenza specialistica appropriata.

\textbf{Risposta.} L'indicazione specialistica sopravvive come concern anche cambiando la scala di priorita' e apre una possibile famiglia autonoma di Decision (tipo di routing, criterio di selezione, confine con servizi sanitari esterni). Forzarla sotto MR-01 allargherebbe retroattivamente lo scope dell'MR.

\textbf{Esito: FAIL come Decision child di MR-01.}

\subsection{Responsibility-boundary / semantic sufficiency}
\begin{ddtacheck}
\textbf{Domanda.} DermaTriage raccomanda una categoria specialistica oppure possiede l'effettiva assegnazione/presa in carico nel sistema sanitario?
\end{ddtacheck}

La fonte mostra un campo/output \texttt{specialist}, ma non documenta la regola di selezione ne' una responsabilita' di scheduling/assignment clinico effettivo.\par
\textbf{Smallest critical proposition:} recommendation/indication vs actual routing/assignment.\par
\textbf{Disposition: NOT SPECIFIED per l'autorita' di routing effettivo.}

\begin{ddtacore}[title={Specialist candidate disposition}]
\textbf{OUTSIDE CURRENT MR-01 DECISION FAMILY.}\par
Non assegnare un identificatore DEC sotto MR-01. Riaprire la \textbf{Macro Project Map} e sottoporre a D1 un candidato MR separato, senza assumere ancora che debba necessariamente sopravvivere come MR autonomo.
\end{ddtacore}

% -----------------------------------------------------------------------------
\section{D2 cumulativo - chiusura della Decision family di MR-01}
\closedtag

Il gate cumulativo non chiede se ogni source fact e' diventato Decision. Chiede se i \textbf{commitment rilevanti che restringono MR-01} sono stati identificati e separati, se i boundary necessari sono espliciti e se gli altri significati sono stati instradati senza invenzione.

\begin{longtable}{L{58mm}L{24mm}L{75mm}}
\toprule
\textbf{Domanda cumulativa} & \textbf{Esito} & \textbf{Motivazione} \\
\midrule
\endhead
Le Decision accettate restringono chiaramente MR-01? & PASS & DEC-01 governa la continuita' del triage senza immagine; DEC-02 governa la convenzione P1-P4 per la priorita' operativa. \\
I commitment indipendenti sono separati? & PASS & P-scale, mapping, SLA e specialist sono stati separati; nessuna Decision monolitica conserva la tabella sorgente intera. \\
Ogni Decision ha esattamente MR-01 come semantic owner? & PASS & DEC-01 e DEC-02 appartengono allo scope di valutazione/priorita'; specialist e' stato espulso dal branch. \\
I responsibility boundary necessari sono espliciti? & PASS & DermaTriage possiede il triage e la P-scale; B4 e' consumed service. L'effettivo routing specialistico resta fuori dal branch e non viene attribuito per inferenza. \\
Sono rimaste Decision nascoste nei dettagli tecnici? & PASS con routing & Four-stage pipeline, named models e threshold restano realization/config evidence; HIGH/MEDIUM/LOW non viene promosso a Decision autonoma senza evidence di commitment indipendente. \\
I source gap sono preservati? & PASS & SLA semantics, equivalenza symptom-only/full-pipeline e autorita' diagnostica restano NOT SPECIFIED. \\
Esistono obblighi operativi gia' visibili ma non ancora scritti? & YES, non bloccante D2 & Exact urgency/confidence -> P mapping e' future FR candidate sotto DEC-02. D3 non viene anticipato. \\
Esistono meaning fuori MR-01 che richiedono macro review? & YES & Specialist routing, doctor validation/learning lifecycle e diagnostic-support meaning richiedono Macro Project Map review. \\
La Decision family di MR-01 puo' chiudersi senza inventare significato? & PASS & Gli elementi appartenenti al branch sono stabilizzati o correttamente instradati. \\
\bottomrule
\end{longtable}

\begin{ddtacore}[title={D2 cumulative disposition - MR-01}]
\textbf{PASS - DECISION FAMILY MR-01 CLOSED.}\par
DEC-01 e DEC-02 costituiscono la famiglia Decision stabilizzata nello scope corrente di MR-01. SLA non viene promosso; specialist viene rinviato alla Macro Project Map; realization evidence resta non normativa.\par
\textbf{Decision-layer STOP per MR-01: YES.}
\end{ddtacore}

\subsection{Importante: D2 PASS non significa ancora ``iniziare subito gli FR''}
La chiusura D2 autorizza strutturalmente la successiva decomposizione FR sotto DEC-01/DEC-02. Tuttavia il holdout mantiene aperta una condizione upstream gia' nota da R1: la \textbf{Macro Project Map globale era PENDING}. Il candidate specialist ha ora fornito evidence concreta che esistono altri concern macro potenziali.

Per ridurre il rischio di ownership error e cross-MR contamination, la prosecuzione di ricerca torna quindi temporaneamente ad A2 globale prima di A5:

\begin{lstlisting}
MR-01 D2 family closure   PASS
          |
          v
re-open Macro Project Map / D1 breadth review
          |
          +--> specialist-routing concern
          +--> doctor validation / correction concern
          +--> learning / retraining lifecycle concern
          `--> diagnostic-support boundary concern
          |
          v
only after macro ownership is stable
          -> FR authoring under accepted Decisions
\end{lstlisting}

Questa non e' una regressione del metodo: e' l'applicazione del principio di ownership prima della decomposizione downstream.

% -----------------------------------------------------------------------------
\section{Valutazione metodologica dopo la Decision family di MR-01}

\subsection{Pregi confermati}
\begin{enumerate}
  \item \textbf{Gate visibility in validation mode.} Rendere esplicite le domande ha permesso di osservare dove avviene realmente il judgement e di distinguere la validazione sperimentale dall'authoring operativo futuro.
  \item \textbf{Split gate.} Ha trasformato il candidato monolitico P-scale + mapping + SLA + specialist in significati con owner diversi.
  \item \textbf{Decision-vs-realization.} Ha impedito la promozione automatica di adaptation layer, named models, threshold e pipeline tecnica.
  \item \textbf{Semantic sufficiency.} Ha mostrato che valori numerici precisi possono essere semanticamente incompleti, come per gli SLA.
  \item \textbf{Semantic-owner discipline.} Il candidate specialist ha riaperto il layer MR invece di essere forzato sotto MR-01 per comodita' gerarchica.
  \item \textbf{No layer filling.} Un candidate Decision puo' terminare in NOT SPECIFIED o OUTSIDE senza che il metamodel venga riempito artificialmente.
\end{enumerate}

\subsection{Costi / pressure da distinguere}
\begin{itemize}
  \item \textbf{Methodological cognitive load: reale.} Split, semantic owner e Decision-vs-FR/realization richiedono judgement disciplinato.
  \item \textbf{Validation-record verbosity: intenzionale.} In questa fase scriviamo domanda e risposta per studiare il metodo; non va confusa con un requisito del documento finale.
  \item \textbf{Operational authoring overhead: ancora da misurare.} Un editor puo' rendere transienti molte domande e persistere solo il significato governato e le eccezioni materialmente rilevanti.
  \item \textbf{Breadth/depth sequencing pressure.} Il candidate specialist dimostra che una review troppo depth-first puo' scoprire tardi un owner macro diverso. La guida raccomanda breadth-first; il holdout dovra' verificare se serve rendere piu' forte questa raccomandazione.
\end{itemize}

\subsection{Observation register aggiornato}
\begin{longtable}{L{45mm}L{112mm}}
\toprule
\textbf{Observation} & \textbf{Interpretazione corrente} \\
\midrule
\endhead
OBS-DDTA-DERMA-01 - Gate visibility & In validation mode question + evidence + answer + disposition devono essere visibili; non ancora requisito di persistenza nell'authoring operativo. \\
OBS-DDTA-DERMA-02 - Validation vs assisted authoring & Le domande possono diventare checklist/prompt transienti in un editor maturo; la ricerca conserva invece le risposte per audit metodologico. \\
OBS-DDTA-DERMA-03 - Cross-layer classification assistance & Tooling futuro dovrebbe enfatizzare split, Decision-vs-FR e Decision-vs-realization, senza automatizzare il semantic owner. \\
OBS-DDTA-DERMA-04 - Numeric precision != semantic sufficiency & Threshold, durata, percentuale o SLA devono avere semantic binding esplicito prima di diventare obblighi governati. \\
OBS-DDTA-DERMA-05 - Ownership can force upstream reopen & Un candidate downstream puo' dimostrare che il macro owner non e' ancora stabile; il processo deve permettere un reopen controllato senza retrofitting gerarchico. \\
\bottomrule
\end{longtable}

\begin{ddtacore}[title={Valutazione provvisoria della metodologia dopo D2 MR-01}]
Non emerge un difetto strutturale che richieda una modifica immediata di R4. I gate disponibili hanno prodotto rework e routing coerenti. Il principale punto da verificare e' ora il sequencing breadth-first: la guida lo raccomanda gia', ma il holdout mostrera' se l'operational authoring necessita di una guardrail piu' esplicita prima degli FR.
\end{ddtacore}


\section{Protocollo guidato riutilizzabile per i prossimi elementi}

Questa sezione definisce il formato con cui il documento dovra' essere esteso durante la continuazione. L'obiettivo e' che ogni decisione documentale resti auditabile.

\subsection{Worksheet per ogni candidate DDTA element}

\begin{ddtacheck}[title={1. Source evidence}]
Quali frasi, tabelle, diagrammi o altri artifact del package sostengono il significato? L'evidence descrive project responsibility/commitment oppure soltanto current realization, configuration, runtime o test state?
\end{ddtacheck}

\begin{ddtacheck}[title={2. Candidate project meaning}]
Qual e' la minima proposizione governata che vogliamo preservare? E' scritta con vocabolario di progetto e al livello di astrazione pertinente?
\end{ddtacheck}

\begin{ddtacheck}[title={3. Semantic owner}]
Il significato appartiene a framing, \MR, \DEC, \FR, \SR, \SECR, boundary cross-MR, oppure e' supporting realization/evidence senza collocazione normativa automatica?
\end{ddtacheck}

\begin{ddtacheck}[title={4. Relevant gates}]
Applicare tutti i gate pertinenti, non soltanto quello del tipo candidato. Per una Decision, ad esempio, includere responsibility-boundary e Decision-vs-implementation evidence. Per un FR includere parentage, coherent-unit/split, binding e failure semantics quando materiali.
\end{ddtacheck}

\begin{ddtacheck}[title={5. Explicit disposition}]
Registrare una delle disposition operative: \texttt{PASS}, \texttt{REWORK}, \texttt{NOT SPECIFIED}, \texttt{CONFLICTING}, \texttt{OUTSIDE CURRENT DDTA DOCUMENTATION}, \texttt{METHOD PRESSURE}.
\end{ddtacheck}

\begin{ddtacheck}[title={6. Stopping decision}]
Se il significato corrente e' sufficiente, fermare il branch. Non creare elementi inferiori per completare visivamente una gerarchia.
\end{ddtacheck}

\subsection{Ordine rimanente della review}

\begin{longtable}{L{17mm}L{39mm}L{90mm}}
\toprule
\textbf{Step} & \textbf{Review} & \textbf{Domanda guida} \\
\midrule
\endhead
A3 & Semantic sufficiency & Sopravvivono letture materialmente diverse? Qual e' la smallest critical proposition? \\
A4 & Decision & Quale commitment materiale restringe l'MR? E' governato o solo implementation evidence? \\
A5 & FunctionalRequirement & Quale obbligo operativo independently assessable discende dalla Decision? \\
A6 & Requirement split & Una clausola puo' cambiare o essere assessed indipendentemente dalle altre? \\
A7 & SpecializedRequirement & Esiste una proprieta' concern-specific aggiuntiva che strengthening il parent FR? \\
A8 & SecurityRequirement & Esiste una singola security property governata con failure mode identificabile? \\
A9 & Cross-MR / consumed service & Il branch possiede o consuma la capability? Quale garanzia del servizio e' davvero governata? \\
A10 & Terminology / bindings & Le stesse identita' usano termini e riferimenti stabili? \\
A11 & Propagation & Le correzioni upstream sono preservate nei discendenti senza narrowing/widening silenzioso? \\
A12 & Completeness / promotion & La candidate baseline e' completa, coerente e accettabile come source baseline per BA? \\
\bottomrule
\end{longtable}

% -----------------------------------------------------------------------------
\section{Registro delle osservazioni metodologiche}

Le osservazioni in questa sezione sono volutamente separate dalla documentazione di progetto. Non devono contaminare MR/Decision/Requirement finali. Verranno consolidate solo al termine della documentazione e, successivamente, confrontate con la regressione Documentation $\leftrightarrow$ BA.

\subsection{Pregi osservati finora - provvisori}

\begin{enumerate}
  \item \textbf{Authority discipline.} La metodologia impedisce di prendere come verita' il candidate DDTA precedente, il modello corrente o la BA futura. Questo rende piu' chiaro da dove proviene ogni significato.
  \item \textbf{Problem/solution separation.} Il framing costringe a rimuovere tecnologia e architettura quando non sono intrinseche al problema. Il risultato e' una responsabilita' piu' stabile nel tempo.
  \item \textbf{Ambiguity surfacing.} La distinzione triage vs diagnosi definitiva, facilmente nascosta da endpoint e campi chiamati \texttt{diagnose}/\texttt{pathology}, viene resa esplicita come \NS{} invece di essere risolta per supposizione.
  \item \textbf{Controlled decomposition.} Il gate di STOP impedisce di produrre documenti inferiori soltanto per riempire la gerarchia, mentre il parentage mantiene la tracciabilita' quando la decomposizione e' necessaria.
  \item \textbf{Classification before method change.} Un problema incontrato viene prima classificato come source gap, guide problem, metamodel pressure, L2 issue o downstream issue. Questo riduce il rischio di cambiare il metodo per ogni anomalia documentale.
\end{enumerate}

\subsection{Costi, limiti e pressure da verificare - non conclusivi}

\begin{enumerate}
  \item \textbf{Review effort.} L'applicazione esplicita dei gate richiede piu' lavoro di una riscrittura informale della documentazione. Dovremo valutare se il maggior costo produce un beneficio proporzionato in chiarezza e regressione.
  \item \textbf{Judgement burden.} La distinzione fra project commitment e current realization non e' completamente automatizzabile; il metodo richiede judgement governato e quindi disciplina del reviewer.
  \item \textbf{Technical-realization preservation.} Dataset, modelli, threshold, training configuration e test evidence possono essere materialmente utili senza essere Requirements o Decisions. La loro rappresentazione non normativa resta una pressure area aperta nella R4.
  \item \textbf{Validation-record verbosity.} La verbosita' osservata e' in parte intenzionale nella Research / Validation mode: domanda, evidence e risposta vengono persistite per studiare il metodo. Non va trattata automaticamente come overhead dell'authoring operativo, dove un editor puo' presentare molte domande come checklist transienti.
  \item \textbf{Tooling opportunity with authority boundary.} Molti check strutturali possono essere assistiti da tool, ma semantic owner, equivalenza e meaning non devono essere decisi automaticamente senza governed review.
\end{enumerate}

\subsection{Come DDTA sta gia' migliorando la documentazione DermaTriage}

\begin{tabularx}{\textwidth}{L{52mm}Y}
\toprule
\textbf{Problema tipico nella fonte classica} & \textbf{Miglioramento prodotto dalla review DDTA} \\
\midrule
Purpose, architecture e implementation facts sono mescolati. & Il framing separa il problema stabile dalle scelte correnti di soluzione. \\
Terminologia come \texttt{diagnose}, \texttt{pathology} e ``triage'' puo' suggerire authority diverse. & La semantic sufficiency isola la proposizione critica e conserva \NS{} invece di scegliere una lettura. \\
Una pipeline tecnica suggerisce naturalmente una decomposizione per componenti. & L'MR viene definito per responsibility, non per stage architetturale. \\
Metriche, threshold e configuration sembrano ``importanti'' e quindi rischiano di diventare requisiti per default. & La parameter/evidence classification richiede prima di stabilirne il semantic owner. \\
La mancanza di informazione puo' essere nascosta da wording generico. & \NS{} e \texttt{CONFLICTING} diventano esiti espliciti e tracciabili. \\
\bottomrule
\end{tabularx}

\begin{ddtaworking}[title={Criterio per la valutazione finale della metodologia}]
Alla fine del caso DermaTriage non giudicheremo DDTA in base a quanto ``bella'' appare la gerarchia. Valuteremo invece se il metodo ha: preservato il significato source-supported; reso visibili gap e conflitti; ridotto ambiguity e implementation leakage; prodotto una baseline comprensibile da un reviewer di progetto; consentito una BA fedele senza invenzione; richiesto un costo di authoring/review ragionevole rispetto ai benefici.
\end{ddtaworking}

% -----------------------------------------------------------------------------
\section{Forma prevista della documentazione finale DermaTriage}

Questo working record non impone in anticipo quali documenti debbano esistere. La forma finale sara' determinata dai gate. Tuttavia, se i branch richiederanno l'intera decomposizione corrente, la relazione di ownership restera':

\begin{lstlisting}
Project problem framing
    -> MacroRequirement
        -> Decision
            -> FunctionalRequirement
                -> SpecializedRequirement
                    -> SecurityRequirement
\end{lstlisting}

Ogni branch potra' fermarsi prima quando semanticamente sufficiente. La documentazione finale dovra' essere piu' corta e piu' leggibile di questo registro: conterra' project meaning governato, non il processo di ricerca che lo ha prodotto.

\begin{ddtaprinciple}[title={Working record vs project baseline}]
\textbf{Working record:} evidence, dubbi, gate, alternative, diagnostics, method observations.\par
\textbf{Project baseline finale:} solo significato governato, gap espliciti necessari, termini stabili, responsabilita', commitment e obblighi accettati.
\end{ddtaprinciple}

% -----------------------------------------------------------------------------
% -----------------------------------------------------------------------------
\section{Evoluzione da R2}

\begin{tabularx}{\textwidth}{L{42mm}Y}
\toprule
\textbf{Dimensione} & \textbf{Evoluzione R2 -> R3} \\
\midrule
Rimasto invariato & Authority gate, framing, MR-01, DEC-01, DEC-02, candidate SLA non promosso e D2 cumulativo locale di MR-01 restano validi. \\
Aggiunto & Reopen controllato di A2; MR-02, MR-03, MR-04; dependency \texttt{MR-04 dependsOn MR-03}; candidate \texttt{predicted\_pathology / diagnosis}; cumulative Macro Project Map review. \\
Corretto / instradato & Il candidate specialist esce da MR-01 e diventa MR-02. \texttt{predicted\_pathology} non viene promosso a MR diagnostico e viene routed a MR-01 per classificazione downstream. \\
Gate che ha causato il cambiamento & Semantic-owner review, MR split/merge stress, isolated-MR review, semantic sufficiency e D1 cumulativo. \\
Perche' & Evitare cross-MR contamination, hierarchy filling e promozione di output tecnici a responsabilita' macro non supportate dalla fonte. \\
\bottomrule
\end{tabularx}

\begin{ddtacore}[title={Milestone R3}]
\textbf{A2 / Macro Project Map = CLOSED - D1 CUMULATIVE PASS.}

La chiusura A2 non autorizza ancora gli FR globalmente: MR-03 e MR-04 richiedono Decision review e MR-01 ha un reopen D2 limitato per la classificazione di \texttt{predicted\_pathology}.
\end{ddtacore}

\section{Perche' A2 viene riaperto dopo R2}

R2 ha chiuso la Decision family di \texttt{MR-01} con D2 cumulativo PASS. Durante quella chiusura, tuttavia, il campo/output \texttt{specialist} ha fallito il semantic-owner test rispetto a MR-01: la valutazione di urgenza/priorita' e l'indirizzamento verso una destinazione specialistica sono concerns distinti. La metodologia richiede quindi un reopen controllato del minimo owning layer, cioe' A2 / Macro Project Map, prima di iniziare gli FR.

\begin{ddtacore}[title={Stato ereditato da R2}]
\begin{lstlisting}
MR-01  Valutazione di triage del caso dermatologico
  |- DEC-01 Continuita' del triage senza immagine         PASS
  `- DEC-02 Adozione della P-scale                         PASS

SLA candidate                                            NOT SPECIFIED
specialist routing candidate                            REOPEN A2
\end{lstlisting}
\end{ddtacore}

La regola di lavoro resta breadth-first: prima si stabilizza la Macro Project Map globale, poi si ritorna alle Decision dei nuovi MR e infine agli FR.

% -----------------------------------------------------------------------------
\section{MR-02 candidate - indirizzamento specialistico}

\subsection{Source evidence}

L'Architecture Document afferma che l'adaptation layer assegna \texttt{P1-P4, specialist and SLA}. Il Test Environment richiede nel risultato del full pipeline \texttt{P1-P4 + Italian specialist + SLA time}. Questa evidence supporta l'esistenza di una \emph{indicazione specialistica} associata al caso.

Non emerge invece evidence sufficiente di prenotazione, assegnazione a uno specifico medico, trasferimento del caso, accettazione da parte di un servizio o presa in carico effettiva.

\begin{tabularx}{\textwidth}{L{64mm}L{35mm}Y}
\toprule
\textbf{Critical proposition} & \textbf{Classificazione} & \textbf{Nota} \\
\midrule
Il risultato contiene una indicazione specialistica & AFFIRMED & Ripetuta in architecture e test setup. \\
Esiste una regola documentata di scelta dello specialist & NOT SPECIFIED & Nessuna regola governabile ricostruita. \\
Il sistema assegna realmente il paziente a un medico/servizio & NOT SPECIFIED & \texttt{assigns} non basta a stabilire presa in carico. \\
\bottomrule
\end{tabularx}

\subsection{Merge stress rispetto a MR-01}

\begin{ddtacheck}[title={Domanda}]
Unendo priorita' di triage e indirizzamento specialistico nello stesso MR, costringeremmo insieme responsabilita' o future Decision family sostanzialmente differenti?
\end{ddtacheck}

\textbf{Risposta: SI'.} Priorita' e destinazione specialistica possono cambiare indipendentemente. Una riguarda \emph{quanto e' urgente il caso}; l'altra \emph{verso quale dominio specialistico orientarlo}. La co-locazione nello stesso output o componente non implica semantic co-ownership.

\begin{ddtaanti}[title={Merge disposition}]
\textbf{MERGE INTO MR-01 = FAIL.} MR-01 conserva il boundary gia' stabilito alla determinazione dell'urgenza/priorita'.
\end{ddtaanti}

\subsection{Candidate MR-02}

\begin{ddtacore}[title={MR-02 - Indirizzamento specialistico del caso dermatologico}]
\textbf{Intent.} Determinare un'indicazione di destinazione specialistica per il caso dermatologico, al fine di supportarne il successivo instradamento verso l'assistenza appropriata.

\textbf{Context.} DermaTriage supporta il triage di casi dermatologici e include nel risultato un'indicazione specialistica associata al caso.

\textbf{Stakeholders.} Paziente; professionista sanitario coinvolto nella presa in carico/review.

\textbf{Scope.} Determinazione dell'indicazione specialistica associata al caso.

\textbf{Boundary.} Non implica, in assenza di ulteriore evidence, assegnazione di uno specifico professionista, prenotazione, trasferimento del caso o garanzia di presa in carico.
\end{ddtacore}

\subsection{D1 - domande e gate}

\begin{tabularx}{\textwidth}{L{62mm}L{21mm}Y}
\toprule
\textbf{Domanda} & \textbf{Esito} & \textbf{Risposta / evidence} \\
\midrule
Una sola macro-responsabilita'? & PASS & Determinare una indicazione/destinazione specialistica. \\
Problem anchoring? & PASS & Sopravvive eliminando B4, API, modelli e componenti. \\
Architecture/technology resilience? & PASS & Il concern resta valido cambiando interamente la soluzione. \\
E' soltanto un atomic FR? & PASS MR & \texttt{return specialist string} sarebbe FR; determinare la destinazione specialistica e' responsabilita' macro. \\
Split nasce dalla soluzione? & NO & La distinzione \emph{urgency vs destination} sopravvive alla neutralizzazione. \\
Dependency su MR-01 governata? & NOT SPECIFIED & La fonte non dice che la destinazione dipenda necessariamente dalla priorita'. \\
\bottomrule
\end{tabularx}

\subsection{Semantic sufficiency e wording correction}

La parola sorgente \texttt{assigns} ammette almeno due letture materialmente diverse: raccomandazione di categoria specialistica oppure assegnazione operativa a un professionista/servizio. L'evidence disponibile supporta solo la prima in modo stabile.

\begin{ddtacheck}[title={Smallest critical proposition}]
DermaTriage possiede la responsabilita' dell'effettiva assegnazione/presa in carico specialistica oppure produce soltanto un'indicazione destinata a supportare l'instradamento?
\end{ddtacheck}

\textbf{Indicazione specialistica: AFFIRMED.}\quad
\textbf{Presa in carico/assegnazione effettiva: NOT SPECIFIED.}

La formulazione viene quindi corretta da \emph{assegnazione dello specialista} a \emph{indirizzamento specialistico / indicazione specialistica}.

\begin{ddtacore}[title={MR-02 disposition}]
\textbf{D1 = PASS dopo semantic rewording.}
\end{ddtacore}

\subsection{STOP branch}

\begin{ddtacheck}[title={Domanda STOP}]
La fonte governa ulteriore significato sufficiente per creare almeno una Decision sotto MR-02 senza inventare policy, allowed domain, routing rule o boundary esterni?
\end{ddtacheck}

\textbf{Risposta: NO.} La fonte afferma l'output specialist ma non governa abbastanza la policy che lo determina.

\begin{ddtaprinciple}[title={STOP result}]
\textbf{MR-02 PASS -> STOP BRANCH AT MR.} Nessuna Decision wrapper viene creata per riempire artificialmente la gerarchia.
\end{ddtaprinciple}

% -----------------------------------------------------------------------------
\section{Macro-concern doctor review / adaptation - split iniziale}

La documentazione contiene una catena consistente: doctor validation/correction, raccolta di correzioni, prompt evolution, classifier retraining, quality gates, deployment/rollback. Trattarla come un solo MR \emph{continuous learning} nasconderebbe due responsabilita' autonome.

\begin{ddtacheck}[title={Split question}]
Validazione clinica degli esiti e adattamento futuro del sistema possono esistere e cambiare indipendentemente, con conseguenze e future Decision family differenti?
\end{ddtacheck}

\textbf{Risposta: SI'.} La validazione puo' esistere senza retraining; la policy di adattamento puo' cambiare mantenendo invariata la raccolta della review clinica.

\begin{ddtacore}[title={Split disposition}]
\textbf{PASS -> due candidate macro-responsibilities: MR-03 e MR-04.}
\end{ddtacore}

% -----------------------------------------------------------------------------
\section{MR-03 - gestione della validazione clinica degli esiti}

\subsection{Candidate meaning e responsibility boundary}

Una formulazione iniziale \emph{DermaTriage valida clinicamente i propri risultati} sarebbe errata: il giudizio clinico appartiene al professionista sanitario. La responsabilita' di progetto e' supportare la registrazione e l'uso governato del risultato di quella review.

\begin{ddtacore}[title={MR-03 - Gestione della validazione clinica degli esiti}]
\textbf{Intent.} Supportare la registrazione e l'utilizzo della validazione o correzione effettuata da un professionista sanitario sugli esiti prodotti da DermaTriage, mantenendo distinguibile il risultato della revisione dall'esito originario.

\textbf{Context.} Gli esiti prodotti da DermaTriage possono essere sottoposti a revisione clinica e la documentazione prevede la registrazione di conferme o correzioni del medico.

\textbf{Stakeholders.} Professionista sanitario; soggetti responsabili della governance e dell'evoluzione del sistema.

\textbf{Scope.} Gestione del risultato della revisione clinica associata a un esito DermaTriage.
\end{ddtacore}

\subsection{D1 - domande e gate}

\begin{tabularx}{\textwidth}{L{61mm}L{21mm}Y}
\toprule
\textbf{Domanda} & \textbf{Esito} & \textbf{Risposta / evidence} \\
\midrule
Una sola macro-responsabilita'? & PASS & Governare il risultato della review/correzione clinica. \\
Problem anchoring? & PASS & Resta visibile eliminando endpoint, B4 e storage. \\
Architecture resilience? & PASS & La responsabilita' non dipende da tecnologia corrente. \\
Atomic FR? & PASS MR & La famiglia futura puo' includere chi valida, cosa corregge, correlation e lifecycle dell'esito. \\
Responsibility boundary chiaro? & PASS dopo REWORK & Il medico possiede il giudizio clinico; DermaTriage gestisce il risultato della review. \\
\bottomrule
\end{tabularx}

\subsection{Semantic sufficiency}

\begin{ddtacheck}[title={Smallest critical proposition}]
La correzione del medico sostituisce semanticamente l'esito AI oppure costituisce una validazione/correzione separata e correlata?
\end{ddtacheck}

\begin{tabularx}{\textwidth}{L{70mm}L{33mm}}
\toprule
\textbf{Proposition} & \textbf{Classification} \\
\midrule
Doctor review/correction exists & AFFIRMED \\
Review result is associated with an AI outcome & AFFIRMED \\
Correction overwrites/replaces original outcome & NOT SPECIFIED \\
Clinical correction constitutes definitive diagnosis & NOT SPECIFIED \\
\bottomrule
\end{tabularx}

MR-03 evita entrambe le letture non governate e resta semanticamente stabile.

\begin{ddtacore}[title={MR-03 disposition}]
\textbf{D1 = PASS. STOP = NO.} La fonte contiene ulteriore significato governato downstream sulla review/correction, da analizzare dopo la stabilizzazione breadth-first della Macro Project Map.
\end{ddtacore}

% -----------------------------------------------------------------------------
\section{MR-04 - adattamento controllato sulla base della revisione clinica}

\subsection{Neutralizzazione della soluzione}

Una formulazione \emph{CNN retraining + prompt evolution} sarebbe solution-oriented. Eliminando nomi e meccanismi resta una responsabilita' stabile: modificare controllatamente il comportamento futuro usando evidence derivata dalla revisione clinica, mantenendo criteri per evitare l'accettazione di regressioni.

\begin{ddtacore}[title={MR-04 - Adattamento controllato sulla base della revisione clinica}]
\textbf{Intent.} Consentire l'evoluzione controllata del comportamento di DermaTriage utilizzando evidence derivata dalle validazioni e correzioni cliniche accumulate, evitando che modifiche degradanti vengano accettate senza controllo.

\textbf{Context.} La documentazione prevede processi di aggiornamento del comportamento del sistema attivati da correzioni cliniche e sottoposti a criteri di valutazione e recupero.

\textbf{Stakeholders.} Professionisti sanitari che producono la review; responsabili tecnici/governance del sistema.

\textbf{Scope.} Evoluzione controllata del comportamento del sistema basata su evidence proveniente dalla revisione clinica.
\end{ddtacore}

\subsection{D1 - domande e gate}

\begin{tabularx}{\textwidth}{L{62mm}L{21mm}Y}
\toprule
\textbf{Domanda} & \textbf{Esito} & \textbf{Risposta / evidence} \\
\midrule
Una sola macro-responsabilita'? & PASS & Adattamento governato del comportamento futuro. \\
Problem anchoring? & PASS & Sopravvive a rimozione di model name, prompt, framework e threshold. \\
Architecture resilience? & PASS & Puo' essere realizzato con tecnologie diverse. \\
Temporal stability? & PASS & Trigger e meccanismi possono cambiare senza eliminare la responsabilita'. \\
Prompt evolution e classifier retraining richiedono due MR? & NO SPLIT & Sono meccanismi indipendenti ma la distinzione non sopravvive come due responsabilita' macro neutrali. \\
Quality gates/rollback richiedono MR separato? & NO per ora & L'evidence li lega al lifecycle dell'adattamento; classificazione downstream ancora aperta. \\
\bottomrule
\end{tabularx}

\subsection{Responsibility boundary e semantic sufficiency}

La fonte documenta trigger, quality gate e rollback, ma non chiarisce completamente se il superamento dei gate autorizzi automaticamente il deployment oppure se intervenga un'approvazione esterna.

\begin{ddtacheck}[title={Smallest critical proposition}]
Il superamento dei quality gate determina automaticamente il deployment oppure richiede una decisione/approvazione esterna?
\end{ddtacheck}

\textbf{Governed adaptation mechanisms exist: AFFIRMED.}\quad
\textbf{Autonomous production deployment authority: NOT SPECIFIED.}

MR-04 resta valido perche' non afferma autonomia di rilascio.

\begin{ddtacore}[title={MR-04 disposition}]
\textbf{D1 = PASS. STOP = NO.} Trigger, strategie di adattamento, evidence usata, quality gates e rollback forniscono significato downstream sufficiente per future Decision review.
\end{ddtacore}

% -----------------------------------------------------------------------------
\section{Dependency review: MR-04 dependsOn MR-03}

MR-04 e' definito precisamente come adattamento basato su validazioni/correzioni cliniche accumulate. La relazione quindi non deriva dalla vicinanza architetturale: MR-04 consuma evidence del tipo governato da MR-03.

\begin{ddtacheck}[title={Dependency question}]
Eliminando la disponibilita' di evidence clinica governata da MR-03, il significato documentato di MR-04 conserva ancora il proprio input essenziale?
\end{ddtacheck}

\textbf{Risposta: NO.} Il dependency e' genuino.

\begin{ddtacore}[title={Dependency disposition}]
\texttt{dependsOn(MR-04, MR-03) = PASS}. La relazione non implica containment e non assegna a MR-03 lo storage tecnico consumato da MR-04.
\end{ddtacore}

% -----------------------------------------------------------------------------
\section{Macro Project Map corrente}

\begin{ddtacore}[title={State after MR-02 / MR-03 / MR-04 review}]
\begin{lstlisting}
DermaTriage
|
+-- MR-01  Valutazione di triage del caso dermatologico
|     +-- DEC-01  Continuita' del triage senza immagine
|     `-- DEC-02  Adozione della P-scale
|
+-- MR-02  Indirizzamento specialistico del caso dermatologico
|     `-- STOP AT MR
|
+-- MR-03  Gestione della validazione clinica degli esiti
|     `-- NON STOP
|
`-- MR-04  Adattamento controllato sulla base della revisione clinica
      `-- NON STOP
          dependsOn MR-03
\end{lstlisting}
\end{ddtacore}

Non viene introdotta alcuna relazione \texttt{dependsOn} tra MR-01 e MR-02, perche' la fonte non stabilisce che l'indicazione specialistica debba essere derivata necessariamente dalla priorita'.


% -----------------------------------------------------------------------------
\section{Candidate macro: predicted pathology / diagnosis}

\subsection{Source evidence e tensione semantica}

La documentazione afferma contemporaneamente che DermaTriage e' orientato al \emph{triage dermatologico} e che la pipeline produce informazioni cliniche ulteriori rispetto alla sola urgenza: descrizione clinica, casi simili e \texttt{predicted\_pathology}. Esiste inoltre un endpoint denominato \texttt{/diagnose}.

\begin{tabularx}{\textwidth}{L{67mm}L{33mm}}
\toprule
\textbf{Evidence / proposition} & \textbf{Classificazione} \\
\midrule
Purpose: AI-assisted dermatology triage & AFFIRMED \\
Le fasi arricchiscono il contesto prima della final triage decision & AFFIRMED \\
Stage 2 produce clinical description & AFFIRMED \\
Stage 3 recupera casi simili & AFFIRMED \\
Stage 4 produce \texttt{predicted\_pathology} & AFFIRMED \\
Esiste un endpoint \texttt{/diagnose} & AFFIRMED \\
DermaTriage possiede autorita' di diagnosi clinica definitiva & NOT SPECIFIED \\
\texttt{predicted\_pathology} costituisce capability autonoma dal triage & NOT SPECIFIED \\
\bottomrule
\end{tabularx}

\subsection{Candidate MR-05 e one-responsibility test}

Una formulazione superficiale sarebbe:

\begin{ddtacore}[title={Candidate MR-05 - Predizione patologica del caso dermatologico}]
Determinare la patologia associata al caso dermatologico.
\end{ddtacore}

Astrattamente, produrre una caratterizzazione/ipotesi patologica potrebbe costituire una responsabilita' coerente. Il one-macro-responsibility test da solo quindi non la esclude.

\textbf{One responsibility = PASS, ma D1 non e' ancora chiuso.}

\subsection{Problem anchoring}

\begin{ddtacheck}[title={Domanda}]
Se eliminiamo l'implementazione corrente, il problema generale affrontato dal progetto richiede esplicitamente una responsabilita' autonoma di diagnosi/predizione patologica?
\end{ddtacheck}

\textbf{Risposta: NON STABILITO.} Il framing e la fonte autoritativa presentano il progetto come supporto al triage e descrivono le fasi analitiche come arricchimento del contesto prima della decisione di triage. Non esiste evidence equivalente a una responsabilita' autonoma di diagnosi clinica.

\begin{ddtaopen}[title={Problem-anchoring disposition}]
\textbf{NOT ESTABLISHED.} Non promuovere il candidate MR sulla sola presenza di un output patologico.
\end{ddtaopen}

\subsection{Architecture-neutralization test}

Eliminando \texttt{BioMistral}, \texttt{Qwen2-VL}, \texttt{ChromaDB}, \texttt{/diagnose}, JSON e la numerazione delle stage, resta comunque un significato reale:

\begin{quote}
DermaTriage produce una caratterizzazione/ipotesi patologica come parte dell'analisi del caso.
\end{quote}

Quindi il concern non e' un puro nome di componente. Tuttavia la neutralizzazione non decide ancora se questo risultato sia una macro-responsabilita' autonoma oppure un risultato analitico di supporto al triage.

\textbf{Architecture resilience = PASS; independent MR status = ancora aperto.}

\subsection{Isolated-MR review}

\begin{ddtacheck}[title={Domanda decisiva}]
Rimuovendo MR-05 scompare una vera macro-capability/responsabilita' governata oppure scompare soprattutto un risultato analitico interno a una capability gia' esistente?
\end{ddtacheck}

\textbf{Risposta: la fonte non dimostra una capability autonoma.} Rimuovendo il candidate MR scompare certamente \texttt{predicted\_pathology}, ma il progetto resta concettualmente comprensibile come analisi del caso finalizzata a urgenza e priorita' di triage.

\begin{ddtaopen}[title={Isolated-MR disposition}]
\textbf{FAIL / insufficient evidence for independent MR.}
\end{ddtaopen}

\subsection{Merge stress rispetto a MR-01}

Il fallimento come MR autonomo non autorizza ad ampliare MR-01 fino a \emph{triage + diagnosi}. Una tale fusione attribuirebbe a DermaTriage una responsabilita' diagnostica non stabilita dalla fonte.

\begin{ddtacore}[title={Merge disposition}]
\textbf{NO MR-05; NO expansion of MR-01 to definitive diagnosis.}
\end{ddtacore}

\subsection{Semantic sufficiency}

Due letture restano materialmente differenti:

\begin{enumerate}
  \item \texttt{predicted\_pathology} e' un risultato analitico di supporto al triage;
  \item \texttt{predicted\_pathology} e' un risultato diagnostico autonomamente governato e destinato all'uso clinico.
\end{enumerate}

\begin{ddtacheck}[title={Smallest critical proposition}]
DermaTriage deve produrre una predizione patologica come risultato autonomo destinato all'uso clinico oppure tale predizione e' informazione analitica usata/supportata nel processo di triage?
\end{ddtacheck}

\textbf{Classificazione: NOT SPECIFIED.} Rimane inoltre \textbf{definitive clinical diagnostic authority = NOT SPECIFIED}.

\subsection{Identifier tecnico non equivale a project meaning}

\begin{ddtaanti}
Non inferire \texttt{/diagnose -> system diagnoses -> diagnosis MR}. Un identifier tecnico o un campo \texttt{predicted\_pathology} prova l'esistenza dell'informazione, non la responsabilita' macro che il progetto le attribuisce.
\end{ddtaanti}

\subsection{Disposition e routing}

\begin{ddtacore}[title={Candidate MR-05 disposition}]
\textbf{MR-05 candidate -> NOT SPECIFIED -> NOT PROMOTED.}

L'evidence materiale non viene eliminata. Viene instradata al semantic owner macro piu' coerente con la fonte corrente: \textbf{MR-01 - Valutazione di triage del caso dermatologico}, per classificazione downstream Decision/FR/realization.
\end{ddtacore}

Questo routing non afferma ancora che \texttt{predicted\_pathology} sia un requisito governato di MR-01. Significa soltanto che la successiva domanda corretta e' di livello D2/D3:

\begin{lstlisting}
predicted pathology
        |
        v
governed project commitment?
        |
   +----+----+
   |         |
  YES       NO
   |         |
Decision/FR  realization / supporting evidence
\end{lstlisting}

Il routing puo' quindi richiedere un \textbf{reopen limitato di D2 di MR-01}, senza modificare MR-01 stesso.

% -----------------------------------------------------------------------------
\section{Macro Project Map cumulative review}

\subsection{Copertura e macro story}

La mappa candidata contiene quattro responsabilita' macro stabilizzate:

\begin{tabularx}{\textwidth}{L{24mm}L{64mm}Y}
\toprule
\textbf{MR} & \textbf{Responsabilita'} & \textbf{Stato} \\
\midrule
MR-01 & Valutazione di triage del caso dermatologico & D1 PASS; Decision family locale gia' chiusa in R2, con possibile D2 reopen limitato per pathology routing \\
MR-02 & Indirizzamento specialistico del caso dermatologico & D1 PASS; STOP AT MR \\
MR-03 & Gestione della validazione clinica degli esiti & D1 PASS; NON STOP \\
MR-04 & Adattamento controllato sulla base della revisione clinica & D1 PASS; NON STOP; dependsOn MR-03 \\
\bottomrule
\end{tabularx}

La macro story e' leggibile senza conoscere l'architettura corrente: valutare il caso e la sua priorita', indicarne una destinazione specialistica, governare la review clinica degli esiti e usare tale evidence per un adattamento controllato del comportamento futuro.

\subsection{Overlap review}

\begin{tabularx}{\textwidth}{L{28mm}Y}
\toprule
\textbf{MR} & \textbf{Domanda di dominio dominante} \\
\midrule
MR-01 & Quanto e' urgente il caso / quale priorita' operativa deve risultare? \\
MR-02 & Verso quale assistenza specialistica deve essere orientato il caso? \\
MR-03 & Come viene governato il risultato della revisione/correzione clinica di un esito? \\
MR-04 & Come evolve controllatamente il comportamento futuro usando evidence di review clinica? \\
\bottomrule
\end{tabularx}

Le responsabilita' sono distinguibili e non richiedono fusione.

\subsection{Architecture resilience e dependency}

Nessun MR dipende semanticamente dai nomi \texttt{EfficientNet}, \texttt{Qwen}, \texttt{BioMistral}, \texttt{ChromaDB}, \texttt{B4} o dagli endpoint correnti.

L'unica dependency macro sufficientemente governata resta:

\begin{center}
\texttt{MR-04 dependsOn MR-03}
\end{center}

Non vengono inventate dependency tra MR-01 e MR-02.

\subsection{Material unowned evidence}

\begin{tabularx}{\textwidth}{L{48mm}L{28mm}Y}
\toprule
\textbf{Meaning} & \textbf{Disposition} & \textbf{Routing / nota} \\
\midrule
SLA associati a P1-P4 & NOT SPECIFIED & Semantica temporale insufficiente; non promossi a Decision \\
Definitive clinical diagnostic authority & NOT SPECIFIED & Non inferire diagnosi definitiva \\
\texttt{predicted\_pathology} come responsibility autonoma & NOT SPECIFIED & Candidate MR rifiutato; route a MR-01 downstream classification \\
\bottomrule
\end{tabularx}

Nessun meaning materiale viene quindi semplicemente scartato.

\subsection{D1 cumulative gate}

\begin{ddtacheck}[title={Domanda cumulativa}]
L'insieme MR racconta una macro project story comprensibile, stabile e non dipendente dall'architettura, senza responsabilita' materialmente duplicate o artificialmente create?
\end{ddtacheck}

\begin{ddtacore}[title={Macro Project Map disposition}]
\textbf{SI'. MACRO PROJECT MAP - D1 CUMULATIVE PASS.}

A2 puo' essere chiuso. I candidate meaning non promossi restano esplicitamente preservati e instradati.
\end{ddtacore}

\subsection{Macro Project Map stabilizzata}

\begin{ddtacore}[title={State after cumulative D1 closure}]
\begin{lstlisting}
DermaTriage
|
+-- MR-01  Valutazione di triage del caso dermatologico
|     +-- DEC-01  Continuita' del triage senza immagine
|     `-- DEC-02  Adozione della P-scale
|
+-- MR-02  Indirizzamento specialistico del caso dermatologico
|     `-- STOP AT MR
|
+-- MR-03  Gestione della validazione clinica degli esiti
|     `-- NON STOP
|
`-- MR-04  Adattamento controllato sulla base della revisione clinica
      `-- NON STOP
          dependsOn MR-03

Preserved material meaning:
- SLA semantics                         NOT SPECIFIED
- definitive diagnostic authority       NOT SPECIFIED
- predicted_pathology independent MR    NOT SPECIFIED
  -> route to MR-01 downstream review
\end{lstlisting}
\end{ddtacore}

% -----------------------------------------------------------------------------
\section{Valutazione metodologica cumulativa aggiornata}

\subsection{Pregi osservati}

\begin{itemize}
  \item \textbf{Upstream reopen controllato.} Il semantic-owner failure puo' riaprire il minimo owning layer senza destabilizzare il resto della gerarchia.
  \item \textbf{STOP senza hierarchy filling.} MR-02 resta valido pur non avendo Decision documentabili.
  \item \textbf{Architectural co-location != semantic co-ownership.} Output prodotti insieme possono appartenere a responsabilita' differenti.
  \item \textbf{Lifecycle evidence non automaticamente scartata.} Feedback, quality gate e rollback possono rivelare responsabilita' governate dopo solution neutralization.
  \item \textbf{Output != macro responsibility.} \texttt{predicted\_pathology} e \texttt{/diagnose} non bastano a creare un MR diagnostico.
  \item \textbf{Rejected candidate != discarded evidence.} Il candidate MR-05 viene rifiutato ma il meaning materiale viene preservato e routed downstream.
\end{itemize}

\subsection{Costi / difficolta' osservati}

Il costo cognitivo rimane concentrato su semantic ownership e split/merge multi-criterio. In particolare, un output osservabile puo' sembrare una responsabilita' autonoma anche quando la fonte lo documenta soltanto come parte di un processo piu' ampio. La validation-record verbosity resta invece un costo intenzionale dello studio, non ancora un difetto dell'authoring operativo.

\subsection{Osservazioni di holdout}

\begin{tabularx}{\textwidth}{L{42mm}Y}
\toprule
\textbf{Observation} & \textbf{Significato} \\
\midrule
OBS-DDTA-DERMA-05 & Ownership can force upstream reopen. \\
OBS-DDTA-DERMA-06 & Architectural co-location does not imply semantic co-ownership. \\
OBS-DDTA-DERMA-07 & Independence alone is insufficient for MR split; serve anche solution-neutral responsibility separation. \\
OBS-DDTA-DERMA-08 & Lifecycle evidence may expose governed responsibilities. \\
OBS-DDTA-DERMA-09 & Observable output does not imply MacroRequirement. \\
OBS-DDTA-DERMA-10 & Rejected macro candidates still require routing when they contain material supported meaning. \\
\bottomrule
\end{tabularx}

\begin{ddtaprinciple}[title={Valutazione provvisoria}]
Anche dopo la chiusura cumulativa A2 non emerge ancora la necessita' di modificare R4. Le regole esistenti sono state sufficienti a bloccare sia over-splitting sia falsa promozione di output tecnici. Il holdout continua pero' a mostrare punti in cui un futuro editor dovrebbe rendere visibili domande contestuali su semantic ownership, isolated-MR review e routing dei candidate rifiutati.
\end{ddtaprinciple}

% -----------------------------------------------------------------------------
\section{Protocollo per la prosecuzione}

\begin{ddtacore}[title={Next authorized sequence after A2 closure}]
\begin{lstlisting}
1. A2 / Macro Project Map                         CLOSED - PASS
2. MR-01 limited D2 reopen
   -> classify predicted_pathology downstream
3. MR-03 Decision family
4. MR-04 Decision family
5. cumulative Decision closure across active MR branches
6. only then descend to FunctionalRequirements
7. apply individual D3 + cumulative FR-family review
\end{lstlisting}
\end{ddtacore}

Questo checkpoint R3 chiude A2 / Macro Project Map con D1 cumulativo PASS. Il prossimo lavoro autorizzato e' la completion breadth-first delle Decision: reopen limitato di MR-01 per classificare \texttt{predicted\_pathology}, quindi Decision family di MR-03 e MR-04. Gli FR restano differiti fino alla chiusura cumulativa delle Decision attive.

\appendix
\section{Compact review worksheet}

\begin{longtable}{L{45mm}L{105mm}}
\toprule
\textbf{Field} & \textbf{Da compilare per ogni elemento} \\
\midrule
\endhead
Source artifact / location & Documento, sezione/pagina, tabella o artifact che sostiene il significato. \\
Source proposition & Fatto minimo effettivamente supportato. \\
Candidate DDTA type & MR / Decision / FR / SR / SecR / boundary / evidence-only / unresolved. \\
Candidate wording & Testo minimo nel vocabolario del progetto. \\
Semantic owner & Perche' questo livello e non un altro. \\
Relevant gates & Tutti i gate applicati, non solo quello nominale del tipo. \\
Smallest critical proposition & Distinzione minima quando esistono letture materialmente diverse. \\
Evidence status & AFFIRMED / DENIED / NOT SPECIFIED / CONFLICTING. \\
Disposition & PASS / REWORK / NOT SPECIFIED / CONFLICTING / OUTSIDE / METHOD PRESSURE. \\
STOP decision & Fermare il branch o continuare, con motivo. \\
Method observation & Eventuale OBS-DDTA / HYP-METHOD / MIP / DT-METHOD candidate, separata dal project meaning. \\
\bottomrule
\end{longtable}

\section*{Session snapshot}

\begin{lstlisting}
R1 ORIGIN BASELINE
21575a18ccf2feab62b27e3b046fadcd01c90135

R2 CHECKPOINT BASELINE
127c8a1532a193184ce86dd4c95f7547ce3606b9

R3 REPOSITORY BASELINE
4dd52e02cc962aee49d2c4593ff851802338698d

SOURCE PACKAGE
DermaTriage-Docs-20260830T152637Z-1-001.zip
SHA-256
E9ED2C507BEFB95F54A52084687CD1E8798863AE81CF69D09568864D8CBF280E

A0 AUTHORITY GATE                  PASS
A1 PROJECT PROBLEM FRAMING         PASS
DIAGNOSTIC AUTHORITY               NOT SPECIFIED
A2 MACRO PROJECT MAP               D1 CUMULATIVE PASS / CLOSED
MR-01 TRIAGE                       PASS
  DEC-01                           PASS
  DEC-02                           REWORK -> PASS
  SLA CANDIDATE                    NOT SPECIFIED / NOT PROMOTED
  predicted_pathology              ROUTE -> LIMITED D2 REOPEN
MR-02 SPECIALIST ROUTING           PASS -> STOP AT MR
MR-03 CLINICAL VALIDATION          PASS -> NON STOP
MR-04 CONTROLLED ADAPTATION        PASS -> NON STOP
  dependsOn MR-03                  PASS
MR-05 DIAGNOSTIC CANDIDATE         NOT SPECIFIED / NOT PROMOTED
DECISION COMPLETION                NEXT
FR AUTHORING                       DEFERRED
DOCUMENTATION GATE                 OPEN
BASE ANALYSIS                      NOT STARTED / BLOCKED
\end{lstlisting}

\begin{ddtacore}[title={Continuation invariant}]
Le revisioni successive devono preservare gli esiti gia' accettati e aggiungere evidence/gate in modo cumulativo. Un reopen e' ammesso soltanto se motivato da nuova source evidence, semantic-owner conflict o contraddizione rilevata durante la regression.
\end{ddtacore}

\end{document}
